\begin{notebox}
\small\textbf{Purpose of this lecture.} Lecture 3 built an encryption scheme out of MLWE. Today we build \emph{signature} schemes --- a different goal (proving you hold a secret, not hiding a message) that leads to genuinely different designs. We cover \textbf{ML-DSA (Dilithium)}, which reuses MLWE directly inside a technique called Fiat--Shamir-with-Aborts, and \textbf{FN-DSA (Falcon)}, which uses a different lattice (NTRU) and a different technique (GPV trapdoor sampling) entirely. Both sections end by naming the exact step in each algorithm that Lesson 6's physical attacks target --- that mapping is the main deliverable of this lecture.
\end{notebox}

\section{What a Signature Scheme Needs to Do}
\label{sec:L4:what-a-signature-scheme-needs-to-do}

Encryption (Lecture 3) hides information; a signature does the opposite --- it \emph{proves} something publicly, without hiding anything. A signature scheme has three algorithms: $\mathrm{KeyGen}$ (produces a public/secret key pair), $\mathrm{Sign}(sk, \text{message})$ (produces a signature), and $\mathrm{Verify}(pk, \text{message}, \text{signature})$ (accepts or rejects). The security goal is \textbf{unforgeability}: without $sk$, no one should be able to produce a signature that $\mathrm{Verify}$ accepts, even after seeing many genuine signatures on messages of their choosing.

\begin{rmkbox}{The Core Design Tension}
A signature must depend on the secret key (otherwise anyone could forge it) --- but it must not depend on the secret key in a way that \emph{leaks} it, even after an attacker collects thousands of signatures. Both schemes in this lecture are, at heart, different engineering answers to that one tension. Keep this framing in mind; it's the thread connecting today's material straight through to Lesson 6.
\end{rmkbox}

\section{A Primer on Fiat--Shamir (Needed Before Dilithium Makes Sense)}
\label{sec:L4:a-primer-on-fiat-shamir-needed-before-di}

\begin{defbox}{Fiat--Shamir Transform, Informally}
Start with a three-move \emph{interactive} protocol between a Prover (who knows a secret) and a Verifier: (1) Prover sends a \textbf{commitment} $w$ (built from fresh randomness, independent of the secret); (2) Verifier sends a random \textbf{challenge} $c$; (3) Prover computes a \textbf{response} $z$ that combines $w$'s randomness with the secret and $c$, and the Verifier checks a consistency equation relating $w$, $c$, $z$, and the public key. The Fiat--Shamir transform removes the Verifier: the Prover computes $c$ themself, as a hash of $w$ (and the message being signed), $c = H(w, \text{message})$. Because $H$ behaves unpredictably, the Prover still can't predict $c$ before committing to $w$ --- so the security argument survives, but now there's no interaction: the whole transcript $(w,c,z)$ can be published as a non-interactive signature.
\end{defbox}

\begin{rmkbox}{Why This Detour Was Necessary}
Every part of Dilithium's Sign algorithm below --- the commitment $w$, the challenge $c$, the response $z$ --- is exactly this three-move shape, with the interaction removed via hashing. Without seeing the general pattern first, Dilithium's algorithm looks like an arbitrary pile of steps; with it, each line has an obvious job.
\end{rmkbox}

\section{ML-DSA (Dilithium): Parameters}
\label{sec:L4:ml-dsa-dilithium-parameters}

\begin{defbox}{ML-DSA Parameter Sets}
All parameter sets share the familiar ring $R_q = \mathbb{Z}_q[x]/(x^n+1)$ with $n=256$, this time with $q = 8{,}380{,}417$ (a different prime from Kyber's --- still chosen for NTT-friendliness). The secret is a pair of small vectors $(\vec s_1,\vec s_2) \in R_q^\ell \times R_q^k$; the public key is built from an MLWE-style sample using them.
\begin{center}
\begin{tabular}{|l|c|c|c|}
\hline
\textbf{Parameter set} & $(k,\ell)$ & Signature size & Public-key size \\
\hline
ML-DSA-44 & (4,4) & 2{,}420 bytes & 1{,}312 bytes \\
ML-DSA-65 & (6,5) & 3{,}309 bytes & 1{,}952 bytes \\
ML-DSA-87 & (8,7) & 4{,}627 bytes & 2{,}592 bytes \\
\hline
\end{tabular}
\end{center}
\end{defbox}

\section{ML-DSA: KeyGen, Sign, Verify}
\label{sec:L4:ml-dsa-keygen-sign-verify}

\begin{algorithm}[ht]
\caption{ML-DSA.KeyGen}
\begin{algorithmic}[1]
\State Generate $A \in R_q^{k\times \ell}$ uniformly at random (again, from a short public seed)
\State Sample small $\vec s_1 \in R_q^\ell$, $\vec s_2 \in R_q^k$ (coefficients bounded by a small $\eta$)
\State Compute $\vec t = A\vec s_1 + \vec s_2 \in R_q^k$
\State \Return $pk = (A, \vec t)$, $sk = (\vec s_1, \vec s_2)$
\end{algorithmic}
\end{algorithm}

\begin{rmkbox}{Yes, This Is an MLWE Sample Too}
Line 3 is identical in shape to Kyber's KeyGen (Lecture 3). The two schemes start from the \emph{same} hard problem; what differs entirely is what you do with the key afterward.
\end{rmkbox}

\begin{algorithm}[ht]
\caption{ML-DSA.Sign$(sk = (\vec s_1,\vec s_2),\ \text{msg})$ --- Simplified}
\begin{algorithmic}[1]
\State Sample a masking vector $\vec y \in R_q^\ell$ with somewhat larger small coefficients than $\vec s_1,\vec s_2$ (derived deterministically from $sk$ and \text{msg} via a PRF, not from fresh randomness --- see remark below)
\State Compute the commitment $\vec w = A\vec y \in R_q^k$
\State Compute the challenge $c = H(\vec w, \text{msg})$ --- a specially-shaped ``sparse'' small polynomial
\State Compute the response $\vec z = \vec y + c\vec s_1$
\State \textbf{Rejection check:} if any coefficient of $\vec z$ (or of a related quantity built from $\vec w - c\vec s_2$) is too large in absolute value, \textbf{discard everything and go back to step 1} with fresh $\vec y$
\State Once the check passes, \Return signature $(\vec z, c)$ (in the real scheme, a compact \emph{hint} $\vec h$ is also included, letting the verifier reconstruct $\vec w$'s high-order bits exactly --- omitted here for simplicity)
\end{algorithmic}
\end{algorithm}

\begin{algorithm}[ht]
\caption{ML-DSA.Verify$(pk=(A,\vec t),\ \text{msg},\ (\vec z, c))$ --- Simplified}
\begin{algorithmic}[1]
\State Check that every coefficient of $\vec z$ is within the allowed bound (reject if not)
\State Recompute $\vec w' = A\vec z - c\vec t$
\State Accept if $c = H(\vec w', \text{msg})$; reject otherwise
\end{algorithmic}
\end{algorithm}

\begin{rmkbox}{Why Verification Works: the Algebra}
Substitute: $\vec w' = A\vec z - c\vec t = A(\vec y + c\vec s_1) - c(A\vec s_1+\vec s_2) = A\vec y + cA\vec s_1 - cA\vec s_1 - c\vec s_2 = A\vec y - c\vec s_2 = \vec w - c\vec s_2$. The two copies of $cA\vec s_1$ cancel --- the same style of cancellation as Kyber's decryption in Lecture 3. $\vec w'$ isn't \emph{exactly} $\vec w$; it's $\vec w$ shifted by a small quantity $c\vec s_2$. Real ML-DSA's ``high bits'' / hint mechanism (omitted above) is specifically engineered so that this small shift doesn't change the \emph{high-order} bits of $\vec w$, so $H(\vec w', \text{msg})$ still reproduces the same challenge $c$ the signer used. Our simplified version above sidesteps this by just checking exact equality of the full hash input in a toy example (Section~\ref{sec:L4:a-fully-worked-toy-example}) small enough that no shift occurs.
\end{rmkbox}

Algorithms~1--3 describe three separate computations; Figure~\ref{fig:L4:dsa-protocol} puts them back together to make explicit which party runs which algorithm and what actually crosses the wire --- worth pinning down carefully, since a signature protocol's shape is structurally different from ML-KEM's request/response exchange (Lecture 3, Figure~\ref{fig:L3:cpa-protocol}).

\begin{figure}[ht]
\centering
\begin{tikzpicture}[scale=1, every node/.style={font=\small}]
  \node[font=\bfseries] at (0,0.7) {Signer};
  \node[font=\bfseries] at (7,0.7) {Verifier};
  \draw[dashed, gray] (0,0.3) -- (0,-5.6);
  \draw[dashed, gray] (7,0.3) -- (7,-5.6);

  \node[draw, fill=green!10, minimum width=3.8cm, minimum height=0.9cm, align=center] (kg) at (0,0) {KeyGen\\ $\to (pk,sk)$};

  \draw[-{Latex[length=2mm]}, thick] (0,-1.2) -- (7,-1.2) node[midway, above, font=\tiny] {send $pk=(A,\vec t)$ \textbf{once}};

  \node[draw, fill=blue!10, minimum width=4.0cm, minimum height=0.9cm, align=center] (sign) at (0,-2.4) {Sign$(sk,\text{msg})$\\ $\to (\vec z,c)$};

  \draw[-{Latex[length=2mm]}, thick] (0,-3.6) -- (7,-3.6) node[midway, above, font=\tiny] {send $\text{msg}$ and $(\vec z,c)$};

  \node[draw, fill=green!10, minimum width=4.6cm, minimum height=1.05cm, align=center] (ver) at (7,-4.9) {Verify$(pk,\text{msg},(\vec z,c))$\\ $\to$ accept / reject};
\end{tikzpicture}
\caption{The ML-DSA protocol, party by party. Unlike ML-KEM's exchange (Figure~\ref{fig:L3:cpa-protocol}), where both parties run one algorithm each per session, here only the \textbf{Signer} ever runs KeyGen and Sign --- the secret key $sk=(\vec s_1,\vec s_2)$ never leaves the Signer. $pk=(A,\vec t)$ is generated and published \emph{once} (e.g.\ embedded in a certificate) and then reused for every subsequent message; each new $\text{msg}$ only needs a fresh $(\vec z,c)$ pair traveling alongside it. The \textbf{Verifier} never runs KeyGen or Sign at all --- Verify is the Verifier's only algorithm, and it can be run by anyone holding $pk$, independently, with no reply sent back to the Signer.}
\label{fig:L4:dsa-protocol}
\end{figure}

Notice exactly what \emph{is} and \emph{is not} hidden in this exchange. $\mathrm{Sign}(sk,\text{msg})$ takes the message as an input and returns \emph{only} the signature $(\vec z, c)$ --- that is the standard convention for how a signing algorithm is written in the literature, everywhere, not something specific to Dilithium. It does not mean the message stays secret: Figure~\ref{fig:L4:dsa-protocol}'s ``send $\text{msg}$ and $(\vec z,c)$'' arrow is doing real work, and it has to be, since $\mathrm{Verify}$ needs $\text{msg}$ as an explicit input too (it's hashed together with $\vec w'$ on line 3). A signature is an \emph{authenticator computed about a message}, not a container that hides or transports one --- which is exactly the opposite emphasis from ML-KEM's ciphertext (Lecture 3), where the ciphertext itself \emph{is} the thing being transported and the whole point is that no one without $sk$ can read what's inside it. Carrying the message alongside its signature, in the clear, is the caller's job (the application protocol, not the signature scheme); the signature scheme's only promise is that $(\vec z,c)$ could not have been produced by anyone without $sk$.

\section{Why Rejection Sampling Exists (This Is the Important Part)}
\label{sec:L4:why-rejection-sampling-exists-this-is-th}

Look again at line 4 of Sign: $\vec z = \vec y + c\vec s_1$. This one line is the entire reason Dilithium's signing algorithm has a loop in it (``discard and go back to step 1'') instead of just computing an answer and returning it. This section works through, with small numbers you can check by hand, exactly what goes wrong if that loop is removed, and exactly why putting it back fixes the problem.

\subsection{The Leak: What Happens If Every $\vec z$ Is Published}
\label{sec:L4:the-leak-what-happens-if-every-z-is-publ}

Focus on a single coefficient, so $z = y + cs_1$ is now just one integer equation, with $y$ the masking value, $c$ the (public) challenge, and $s_1$ the (secret) value we're protecting. Suppose $y$ is drawn uniformly at random from a symmetric range, say the $21$ integers $\{-10,-9,\ldots,9,10\}$. Because this range is symmetric around $0$, the \emph{average} value of $y$ over many draws is $0$. But the average value of $z=y+cs_1$ is not $0$ --- it's $cs_1$, since $cs_1$ is added to every single draw:
$$
\mathbb{E}[z] = \mathbb{E}[y] + cs_1 = 0 + cs_1 = cs_1.
$$
This is not just a theoretical concern. Concretely, suppose (unknown to an attacker) $s_1=3$ and, to keep this illustration simple, the challenge happens to be $c=1$ on every signature examined. A concrete run of $20$ signatures might produce a sample average of $z$ around $4.1$; average over $100$ signatures and it settles closer to $3.5$; average over $2{,}000$ signatures and it lands at $3.15$ --- converging, exactly as the formula predicts, toward the true value $cs_1=3$. An attacker does not need to see a single $z$ that ``looks wrong''; they only need \emph{many} $z$'s and a bit of averaging (or, more realistically, since $c$ actually changes with every signature, a short linear regression of the observed $z_i$'s against the known, public $c_i$'s) to recover $s_1$ with growing confidence as more signatures accumulate. This is exactly Lesson 5's ``leak a little, combine a lot'' recipe, aimed here at a signature scheme instead of an encryption scheme.

\subsection{The Fix: Only Publish When It's Provably Safe}
\label{sec:L4:the-fix-only-publish-when-its-provably-s}

The fix is not to make $z$ itself hide $s_1$ better --- there is no way to do that while $z$ still genuinely depends on $s_1$, which it must, or the signature would be forgeable by anyone. The fix instead is to control \emph{which} $z$'s ever get published at all, so that the ones an attacker actually sees carry no statistical trace of $s_1$, even though every individual $z$ still technically depends on it.

Suppose $|cs_1|$ is always bounded by some known constant $\beta$ (in this toy illustration, take $\beta=2$), and instead of publishing every $z$, the signer only publishes $z$ when $|z| \le B$ for a bound $B$ chosen with a safety margin: $B = \gamma_1 - \beta$, where $\gamma_1=10$ is the half-width of $y$'s sampling range. Here that gives $B = 10-2=8$, so only $z$ with $|z|\le 8$ ever get published.

Check what this buys us by working through all three possible shifts $cs_1 \in \{-2,0,2\}$ directly (the true value is always somewhere in $[-\beta,\beta]$; checking the two extremes and the middle is enough to see the pattern):
\begin{itemize}[leftmargin=*]
  \item \textbf{Shift $cs_1=0$.} $z=y$ directly. Accepting $|z|\le 8$ means accepting $y\in\{-8,\ldots,8\}$ --- $17$ equally likely values, each appearing with probability $1/21$ (its original probability as a draw of $y$), restricted to this window.
  \item \textbf{Shift $cs_1=+2$.} $z=y+2$. Accepting $|z|\le8$ means $y\in\{-10,\ldots,6\}$ --- also exactly $17$ values, all comfortably inside $y$'s full range $\{-10,\ldots,10\}$, so every one of them was genuinely reachable. The resulting accepted $z$'s again range exactly over $\{-8,\ldots,8\}$, one value of $y$ per value of $z$.
  \item \textbf{Shift $cs_1=-2$.} $z=y-2$. Accepting $|z|\le8$ means $y\in\{-6,\ldots,10\}$ --- again $17$ values, all inside range, and again the accepted $z$'s land exactly on $\{-8,\ldots,8\}$.
\end{itemize}
In \emph{every} case --- regardless of the actual, unknown value of $cs_1$ anywhere in $[-2,2]$ --- the accepted $z$'s are exactly the $17$ integers $\{-8,\ldots,8\}$, each corresponding to exactly one value of $y$ within its original range. The conditional distribution of $z$, \emph{given that it was accepted}, is identical no matter what $s_1$ actually is. An attacker who only ever sees accepted signatures is looking at a distribution that is provably independent of the secret --- not just hard to distinguish in practice, but mathematically identical across every possible value of the secret consistent with the bound $\beta$.

This is worth contrasting with what happens if the margin is too thin. Suppose instead the bound had been set at $B=9$ (only a $1$-unit buffer, less than $\beta=2$). Now, when $cs_1=-2$, accepting $|z|\le9$ requires $y\in\{-7,\ldots,11\}$ --- but $y=11$ is outside $y$'s actual range $\{-10,\ldots,10\}$, so the value $z=9$ can never occur when $cs_1=-2$. Symmetrically, when $cs_1=+2$, the value $z=-9$ can never occur. The set of \emph{reachable} $z$-values now depends on the sign of the unknown shift $cs_1$ --- a real, exploitable statistical difference between ``$s_1$ tends positive'' and ``$s_1$ tends negative,'' visible simply by noticing which extreme values of $z$ do or do not show up over many signatures. This is exactly why the real bound in Dilithium is set with a margin at least as large as the worst-case size of $c\vec s_1$: not as a rough safety cushion, but as the precise threshold below which the ``every shift gives the identical accepted distribution'' argument above stops holding.

Real ML-DSA runs this same check independently on every coefficient of $\vec z$ (and on a related quantity built from $\vec w - c\vec s_2$), and any coefficient that fails --- on any one of them --- throws away the entire attempt and restarts from a fresh $\vec y$. ``Conditioned on being accepted'' is exactly the phrase for this: not that $z$ itself is secret-independent (it never is), but that the distribution of \emph{published, accepted} $z$'s is.

\subsection{Determinism: Fixing One Failure Mode, Creating Another}
\label{sec:L4:determinism-fixing-one-failure-mode-crea}

Step 1 of Sign notes that $\vec y$ is derived \emph{deterministically} from $(sk,\text{msg})$ via a pseudorandom function, rather than from fresh random bits on every call. This choice mirrors a well-known fix (in the style of RFC 6979) for a classical signature failure mode. If two different messages were ever signed using the \emph{same} $y$ by accident --- for instance, because a broken or poorly-seeded random-number generator repeated itself --- an attacker who sees both signatures can subtract the two resulting equations:
$$
z_1 = y+c_1s_1, \qquad z_2 = y+c_2s_1 \qquad \Longrightarrow \qquad z_1-z_2 = (c_1-c_2)s_1,
$$
and solve directly for $s_1 = (z_1-z_2)(c_1-c_2)^{-1}$, since $c_1,c_2,z_1,z_2$ are all public. This is not a hypothetical: an accidentally repeated nonce in classical ECDSA is exactly the bug behind several real, catastrophic key-recovery incidents, including the Sony PlayStation~3 signing-key compromise. Deriving $y$ deterministically from $(sk,\text{msg})$ closes this specific hole, since honestly signing the \emph{same} message twice now reproduces the same signature rather than silently reusing $y$ across two \emph{different} messages --- there is no longer any accidental way for two different messages to end up sharing a $y$.

But determinism trades one failure mode for another. If an active fault attack --- a voltage glitch, a clock glitch, or a deliberately skipped instruction --- can force the signing device to reuse the same $y$ across two \emph{genuinely different} signing calls (for instance, by corrupting the PRF's internal state, or by replaying an earlier internal value back into the computation), the exact same subtraction attack becomes available again. Nothing about either individual signature looks wrong: both still verify correctly, and the reused $y$ leaves no visible trace in either signature on its own. The only difference from the classical bug is \emph{who} caused the repeat and \emph{why} --- an attacker, deliberately, instead of a broken random-number generator, by accident. This ``fault attacks on determinism'' failure mode is precisely the attack surface Lesson 6 targets in detail (and which the practice exercise at the end of Section~\ref{sec:L4:a-fully-worked-toy-example} already has you work out numerically): determinism is a genuine defense against one failure mode, and a genuine attack surface for a different one, and real published Dilithium fault attacks exploit exactly this tension.

\section{A Fully Worked Toy Example}
\label{sec:L4:a-fully-worked-toy-example}

\begin{exbox}{ML-DSA-Style Signing, $k=\ell=1$, $n=4$, $q=17$ (Verified by Hand)}
Reusing $R_{17}=\mathbb{Z}_{17}[x]/(x^4+1)$ and $A = 3+5x+2x^2+x^3$ from Lecture 3.

\textbf{KeyGen.} Let $s_1 = 1$ (constant polynomial) and $s_2 = 1-x$ (small). Then
$$
t = A\cdot s_1 + s_2 = A + (1-x) = (3+1) + (5-1)x+2x^2+x^3 = 4+4x+2x^2+x^3.
$$

\textbf{Sign.} Mask $y = x$. Commitment: $w = A\cdot y = A\cdot x = 16+3x+5x^2+2x^3$ (computed exactly as in Lecture 2's ring-multiplication example). For this toy example, take the challenge to be the simplest possible sparse polynomial, $c=1$ (constant). Response: $z = y + c\cdot s_1 = x + 1 = 1+x$. (No rejection triggered --- coefficients of $z$ are small.) Signature: $(z,c) = (1+x,\ 1)$.

\textbf{Verify.} Recompute $w' = A\cdot z - c\cdot t$:
\begin{align*}
A\cdot z &= A\cdot(1+x) = A + A\cdot x = (3,5,2,1) + (16,3,5,2) \\
&= (19,8,7,3) \equiv (2,8,7,3) \bmod 17, \\
c\cdot t &= 1\cdot t = (4,4,2,1), \\
w' &= (2-4,\,8-4,\,7-2,\,3-1) = (-2,4,5,2) \equiv (15,4,5,2) \bmod 17.
\end{align*}
Compare to the original $w = (16,3,5,2)$: they are \emph{not} identical --- $w' = w - c\cdot s_2 = w - (1,-1,0,0)$, matching the algebra of Section~\ref{sec:L4:ml-dsa-keygen-sign-verify}'s remark exactly (subtract $s_2=1-x=(1,-1,0,0)$ from $w=(16,3,5,2)$: $(16-1,3-(-1),5,2)=(15,4,5,2)$ $=w'$, confirmed). In this toy example we accept the signature by checking this exact algebraic identity holds (the real scheme's hint mechanism is what lets a verifier who does \emph{not} know $s_2$ still confirm this identity, via the high-order bits of $w$ instead of $w$ itself --- omitted here for tractability, as flagged in Section~\ref{sec:L4:ml-dsa-keygen-sign-verify}).
\end{exbox}

\begin{exercisebox}{Practice}
Using the same $A$, $s_1$, $s_2$: sign the same message again, but this time with $y = 1$ (instead of $y=x$) and the same challenge $c=1$. Compute the new $z$ and confirm $w' = w - c\cdot s_2$ still holds with the new $w = A\cdot y$. Then, pretending you are an attacker who somehow learned that the \emph{same} $y=x$ was used for two different messages with challenges $c_1=1$ and $c_2=2$ respectively, write down the two equations for $z_1$ and $z_2$ and show algebraically how subtracting them isolates $s_1$ --- this is the calculation behind Section~\ref{sec:L4:determinism-fixing-one-failure-mode-crea}'s fault-attack discussion, worked out with real numbers.
\end{exercisebox}

\section{FN-DSA (Falcon): a Different Lattice, a Different Technique}
\label{sec:L4:fn-dsa-falcon-a-different-lattice-a-diff}

Falcon abandons MLWE entirely and uses a different hard problem on a different kind of lattice --- worth seeing, since comparing attack surfaces \emph{across} schemes, not just understanding one in isolation, is exactly what real security analysis requires.

\begin{defbox}{NTRU Lattices}
Fix $R_q = \mathbb{Z}_q[x]/(x^n+1)$ as before. An NTRU key pair consists of two \emph{short} polynomials $f,g \in R_q$ (the secret), from which the public key $h = g\cdot f^{-1} \bmod q$ is computed (assuming $f$ is invertible mod $q$). The \textbf{NTRU lattice} associated to $h$ is
$$
\Lambda_h = \{\, (u,v) \in R_q^2 : u + v\cdot h \equiv 0 \!\!\pmod q \,\}.
$$
$(f,g)$ (suitably arranged) is a \emph{short vector} in $\Lambda_h$ --- in fact, a short enough one to serve as a trapdoor \emph{basis} for the whole lattice, not just a single short vector: this is the ``all vectors short together'' property Lecture 1's optional Minkowski's Second Theorem box was pointing toward.
\end{defbox}

\begin{defbox}{The GPV Framework (Gentry--Peikert--Vaikuntanathan)}
Given a trapdoor (short basis) for a lattice $\Lambda$, and a target point $t = H(\text{msg})$ (hashed into the same space as $\Lambda$), \textbf{sample} a lattice point $v \in \Lambda$ close to $t$, using the trapdoor to guide a randomized search --- specifically, sampling from a discrete Gaussian distribution (Lecture 2, Section~\ref{sec:L2:what-the-error-distribution-actually-is}) centered near $t$. The signature is (essentially) the short offset $s = t - v$. Verification recomputes $t$ from the message, checks $s$ is short, and checks $s+v'$ (reconstructed from the public key and $s$) lands back on a valid target. Unlike Dilithium's commit-challenge-response shape, this is a single sampling step against a target --- no interactive protocol being flattened, just a direct trapdoor-guided search.
\end{defbox}

\begin{rmkbox}{Falcon's KeyGen/Sign/Verify, at a Glance}
\textbf{KeyGen:} generate short $f,g$ (this step itself is delicate number theory --- specific polynomial sampling techniques --- which we do not detail here); compute $h=g f^{-1} \bmod q$; the trapdoor is $(f,g)$ (or an equivalent short basis derived from them), the public key is $h$. \textbf{Sign:} hash the message to a target point, then use the trapdoor with GPV sampling (Falcon's specific method is called \emph{fast Fourier sampling}, chosen for speed) to find a short vector $(s_1,s_2)$ with $s_1 + s_2 h \equiv H(\text{msg}) \pmod q$; output $(s_1,s_2)$ (compressed) as the signature. \textbf{Verify:} recompute $s_1 = H(\text{msg}) - s_2 h \bmod q$ from the public $h$ and the received $s_2$, and check that $(s_1,s_2)$ is short enough to plausibly come from the trapdoor.
\end{rmkbox}

Exactly as with ML-DSA (Figure~\ref{fig:L4:dsa-protocol}), it's worth making explicit which party runs which of Falcon's three algorithms, since the same one-time-key, one-signature-per-message shape holds here too --- only the technique \emph{inside} Sign has changed.

\begin{figure}[ht]
\centering
\begin{tikzpicture}[scale=1, every node/.style={font=\small}]
  \node[font=\bfseries] at (0,0.7) {Signer};
  \node[font=\bfseries] at (7,0.7) {Verifier};
  \draw[dashed, gray] (0,0.3) -- (0,-6.0);
  \draw[dashed, gray] (7,0.3) -- (7,-6.0);

  \node[draw, fill=green!10, minimum width=4.4cm, minimum height=1.05cm, align=center] (kg) at (0,0) {KeyGen: sample short $f,g$\\ $h=g f^{-1}\bmod q$};

  \draw[-{Latex[length=2mm]}, thick] (0,-1.4) -- (7,-1.4) node[midway, above, font=\tiny] {send $pk=h$ \textbf{once}};

  \node[draw, fill=blue!10, minimum width=4.8cm, minimum height=1.15cm, align=center] (sign) at (0,-2.8) {Sign: GPV/FFT-sample $(s_1,s_2)$\\ using trapdoor $(f,g)$};

  \draw[-{Latex[length=2mm]}, thick] (0,-4.2) -- (7,-4.2) node[midway, above, font=\tiny] {send $\text{msg}$ and $(s_1,s_2)$};

  \node[draw, fill=green!10, minimum width=5.0cm, minimum height=1.15cm, align=center] (ver) at (7,-5.6) {Verify: recompute $s_1$ from $h,s_2$;\\ check short \& consistent};
\end{tikzpicture}
\caption{FN-DSA's protocol shape, side by side with Figure~\ref{fig:L4:dsa-protocol}: identical party structure to ML-DSA --- only the \textbf{Signer} ever touches the trapdoor $(f,g)$, and only the Signer runs KeyGen and Sign; the \textbf{Verifier} holds nothing but the public $h$ and runs Verify alone, with no reply sent back. What changes is entirely internal to the Signer's Sign step: instead of Dilithium's commit/challenge/response loop with rejection sampling, Falcon runs a single GPV trapdoor-guided sampling step (fast Fourier sampling) directly against the hashed message target --- the external protocol shape is the same, but there is no ``discard and retry'' loop here the way there is in Dilithium's Sign (Algorithm~2).}
\label{fig:L4:falcon-protocol}
\end{figure}

\begin{previewbox}{Why Falcon's Sampler Is the Star of Lesson 6}
Two things make Falcon's signing step uniquely attack-prone compared to Dilithium's. First, GPV sampling needs to sample from a discrete Gaussian \emph{exactly} (or very close to it) --- an approximate or biased sampler can leak information about the trapdoor $(f,g)$ through the statistical shape of many signatures, even without any single signature looking wrong. Second, Falcon's fast Fourier sampling is implemented using \textbf{floating-point arithmetic}, for speed --- and floating-point operations have historically been far harder to make constant-time (free of secret-dependent timing or power variation) than the simple integer comparisons Dilithium's rejection sampling relies on. Put together: Dilithium's main defense (rejection sampling) is a young technique but algorithmically simple to implement safely; Falcon's main mechanism (Gaussian trapdoor sampling) is a mathematically older and better-understood technique, but a much harder \emph{implementation} target to make side-channel-resistant. This is exactly the gap that produced the published single-trace and sampler-leakage attacks on Falcon that Lesson 6 will walk through.
\end{previewbox}

\section{Standardization Status}
\label{sec:L4:standardization-status}

ML-DSA was finalized as \textbf{FIPS 204} alongside ML-KEM and SLH-DSA in August 2024, and is NIST's recommended general-purpose signature standard today.\footnote{NIST, \emph{Module-Lattice-Based Digital Signature Standard}, FIPS 204, August 2024.} FN-DSA (Falcon), by contrast, remains a \textbf{draft} standard (to become FIPS 206): as of mid-2026 it has not yet been finalized, with publication expected in late 2026 or early 2027, and major certificate authorities have stated they will not deploy it in production until it is.\footnote{NIST Initial Public Draft, FIPS 206, submitted for public review in late 2025.} Worth stating plainly: attacking a finalized federal standard (ML-DSA) and attacking a still-draft one (FN-DSA) carry different framing, even though both are technically live research targets today.

\section{Comparing the Three Schemes So Far}
\label{sec:L4:comparing-the-three-schemes-so-far}

\begin{center}
\renewcommand{\arraystretch}{1.35}
\footnotesize
\begin{tabular}{|p{2.6cm}|p{3.1cm}|p{3.4cm}|p{3.4cm}|}
\hline
& \textbf{ML-KEM (Kyber)} & \textbf{ML-DSA (Dilithium)} & \textbf{FN-DSA (Falcon)} \\
\hline
Hard problem & Module-LWE & Module-LWE & NTRU lattice (SIS/GPV-flavored) \\
Purpose & Key encapsulation & Signature & Signature \\
Core technique & MLWE encryption + FO transform & Fiat--Shamir-with-Aborts & GPV trapdoor sampling \\
Key defense mechanism & Implicit rejection (Decaps) & Rejection sampling & Exact Gaussian sampling \\
Riskiest implementation step & Re-encryption check & $\vec y$ generation / bound check & Floating-point sampler \\
Standard status (mid-2026) & FIPS 203, finalized & FIPS 204, finalized & FIPS 206, draft \\
\hline
\end{tabular}
\end{center}

\begin{previewbox}{Where This Leads: Lessons 5 and 6}
You now have, for all three schemes, both the correct algorithm \emph{and} a named ``riskiest step.'' Lesson 5 opens by formalizing the general distinction between attacking the math (everything in Lectures 1--4) versus attacking an implementation of it, then unpacks ML-KEM's one named riskiest step into two concrete, worked attacks. Lesson 6 completes the picture with the remaining four --- two for ML-DSA, two for FN-DSA --- six risky steps in total across the two lessons. At that point, every attack should read as ``targeting a specific line of an algorithm you can already write out from memory,'' not as an isolated trick.
\end{previewbox}

\section{Exercises}
\label{sec:L4:exercises}

\begin{exercisebox}{Homework Set}
\begin{enumerate}[leftmargin=*]
  \item In your own words, explain the Fiat--Shamir transform to someone who has never seen an interactive proof --- use the commit/challenge/response structure from Section~\ref{sec:L4:a-primer-on-fiat-shamir-needed-before-di}, but invent your own non-cryptographic analogy (e.g.\ a game show, a magic trick) to illustrate why removing the interactive Verifier and replacing it with a hash still works.
  \item Complete the practice exercise in Section~\ref{sec:L4:a-fully-worked-toy-example} (two-signature key-recovery calculation) if you have not already.
  \item Compare Kyber's re-encryption check (Lecture 3, Section~\ref{sec:L3:from-cpa-secure-encryption-to-a-cca-secu}) and Dilithium's rejection sampling (Section~\ref{sec:L4:why-rejection-sampling-exists-this-is-th} above): both exist to prevent a secret-dependent quantity from becoming visible to an attacker. In 4--5 sentences, describe what each mechanism is protecting against, and why the two mechanisms look so different from each other despite that similarity of purpose.
  \item Falcon's signature is a pair $(s_1,s_2)$ with $s_1+s_2h\equiv H(\text{msg})\pmod q$. Suppose $h = 5$ (a toy scalar stand-in, not a real ring element) and $H(\text{msg})=12 \pmod{17}$. If $s_2 = 2$, what must $s_1$ be? Is this pair unique, or could other $(s_1,s_2)$ pairs also satisfy the equation --- and if so, what property (from Section~\ref{sec:L4:fn-dsa-falcon-a-different-lattice-a-diff}) picks out the one Falcon actually outputs?
  \item \textbf{Looking ahead:} using the comparison table in Section~\ref{sec:L4:comparing-the-three-schemes-so-far}, predict --- before Lessons 5--6 reveal the real answer --- which of the three schemes' riskiest steps you would guess is \emph{easiest} to defend against with constant-time coding practices alone, and which would need dedicated hardware countermeasures. Briefly justify your guess; we will revisit it directly there.
\end{enumerate}
\end{exercisebox}
